


%%%%%%%%%%%%%%%%%%%%%%%%%%%%%%%%%%%%%%%%%%%%%%%%%%%%%%%%%%%%%%%%%%%%%
%%%%%%%%%%%%%%%%%%%%%%%%%%%%%%%%%%%%%%%%%%%%%%%%%%%%%%%%%%%%%%%%%%%%%
%%%%%%%%%%%%%%%%%%%%%%%%%%%%%%%%%%%%%%%%%%%%%%%%%%%%%%%%%%%%%%%%%%%%%


\begin{frame}{Conclusion}

\begin{itemize}
\item You may be concerned if you could reverse your conclusion by removing
a small proportion of your data.

\pause \item We can quickly and automatically find an approximate influential
set which is accurate for small sets in well-behaved problems.

\pause \item Our influential points are not a coreset.  (But the influence
function can be used in other ways to help find them.)

\pause \item Data dropping is not a metric on distributions, but it
can define a notion of ``size'' of a distribution shift.

\end{itemize}

\end{frame}




%%%%%%%%%%%%%%%%%%%%%%%%%%%%%%%%%%%%%%%%%%%%%%%%%%%%%%%%%%%%%%%%%%%%%
%%%%%%%%%%%%%%%%%%%%%%%%%%%%%%%%%%%%%%%%%%%%%%%%%%%%%%%%%%%%%%%%%%%%%
%%%%%%%%%%%%%%%%%%%%%%%%%%%%%%%%%%%%%%%%%%%%%%%%%%%%%%%%%%%%%%%%%%%%%

\begin{frame}{Links and references}

\footnotesize

% {\bf A workshop paper: }\newline
Ryan Giordano, Rafe Meager, Tamara Broderick
(Giordano and Meager are joint first authors) \\
``An Automatic Finite-Sample Robustness Metric: Can Dropping a Little Data Change Conclusions?''
\newline {\color{blue}\url{https://arxiv.org/abs/2011.14999}}
\parencite{broderick:2025:automaticpart1,broderick:2025:automaticpart2}

\hrulefill

Select blog posts with more details: \hspace{1em}\url{https://rgiordan.github.io}
\begin{itemize}
\item \href{https://rgiordan.github.io/robustness/2021/09/17/amip_p_hacking.html}
    {Data dropping sensitivity overcomes p-hacking}
\item \href{https://rgiordan.github.io/amip/2021/12/17/reweighted_colinear_note.html}
    {Collinearity in OLS after dropping}
\item \href{https://rgiordan.github.io/amip/2021/12/01/influence_is_sum.html}
    {Influence functions and sums}
\item \href{https://rgiordan.github.io/amip/2021/11/08/bootstrap_influence.html}
    {Connections to the bootstrap}
% \item \href{https://rgiordan.github.io/differentiability/2021/07/20/norms.html}
%     {When a norm is the quantity of interest}
\end{itemize}


\hrulefill

Related software on \texttt{github}:
\begin{itemize}
\item \href{https://github.com/rgiordan/zaminfluence}
{rgiordan/zaminfluence (for \texttt{R})}
\end{itemize}

\hrulefill

\end{frame}





\begin{frame}[allowframebreaks]{References}
\renewcommand*{\bibfont}{\normalfont\tiny}
\printbibliography[heading=none]
\end{frame}
